\section{Related Work}

\subsection{Agent-Based Social Network Simulation}

Research on social networks has gradually evolved from static structural analysis to dynamic process modeling, and agent-based modeling (ABM) has become an important tool for simulating and analyzing dynamic social networks~\citep{li2026abm, gilbert2008agent, wang2025abm}. An ABM framework typically includes four core elements---agents, environment, interaction, and rules---and studies how individual-level behavior patterns and interaction rules give rise to macro-level social structures or phenomena through local mechanisms~\citep{gilbert2008agent, vespignani2012modelling, gaoj2016universal}. Early ABM studies were often combined with epidemic models and opinion dynamics models to simulate epidemic spreading in complex social networks, opinion evolution in fully connected networks, and panic buying in public emergencies~\citep{wu2024epidemic, alatas2017opinion, fan2024panic}. These approaches are effective at characterizing overall diffusion speed, network-topology effects, and macro-level evolutionary patterns. However, as prior studies have noted, they usually rely on limited decision rules and logical functions, making it difficult to represent the fine-grained cognitive process of “seeing--understanding--responding--re-expressing” within a single system, or to capture explanation, reinterpretation, and emotional regulation in real public-opinion dynamics~\citep{li2026abm, gaoc2024survey}.

\subsection{LLM-Driven Multi-Agent Social Simulation}

The emergence of large language models has opened new possibilities for multi-agent social simulation. Generative Agents proposed by Park et al.~\citep{park2023generative} use LLMs to simulate daily conversations, memory retrieval, and long-term planning of human-like roles, thereby pushing social network research from a “rule-driven symbolic world” toward a “semantics-driven cognitive world”~\citep{li2026abm}. LLM-based multi-agent systems provide greater semantic flexibility in role-playing, text understanding, comment generation, and situational inference, allowing researchers to place individual cognition, group interaction, and platform environments into one simulation loop~\citep{gaoc2024survey, park2023generative}. Recent work has extended this line to different scenarios. Mou et al.~\citep{mou2024socialmovement} combined core LLM agents with ABM agents to simulate collective opinion changes in social movements. Yang et al.~\citep{yang2024oasis} proposed OASIS for million-scale social simulation. Gao et al.~\citep{gao2023s3} and Piao et al.~\citep{piao2025agentsociety} developed systems that simulate information diffusion and collective behavior on social platforms at large scales. In rumor and adversarial settings, Liu et al.~\citep{liu2025rumorsphere} introduced RumorSphere for million-scale rumor propagation, and Zhang et al.~\citep{zhang2024trendsim} proposed TrendSim with a time-aware interaction mechanism for topic manipulation under poisoning attacks. Although these studies have advanced both scale and generation quality, their frameworks still mainly center on abnormal diffusion, adversarial manipulation, or task-specific settings. In real public emergencies, official release timing, frequency, stance, and intervention intensity all shape public-opinion trajectories at different stages, yet existing models still provide limited systematic modeling of such official-intervention variables.

\subsection{Reflection Mechanisms in LLM Agents}

Reflection is a key cognitive mechanism that enables agents to update their judgments and strategies based on prior experience. In LLM agents, Park et al.~\citep{park2023generative} already treated reflection as a core module for synthesizing higher-level cognition from accumulated memories. Yao et al.~\citep{yao2023react} proposed ReAct, which interleaves reasoning and acting so that language models can update intermediate judgments through interaction with the environment. Shinn et al.~\citep{shinn2023reflexion} further proposed Reflexion, which combines language feedback, self-evaluation, and retrying, allowing agents to generate verbal reflections from failures and explicitly reuse them in subsequent decisions. In multi-agent simulation, the value of such mechanisms lies not only in improving task performance, but also in enabling agents to dynamically revise their cognitive states, emotional tendencies, and expressive strategies on the basis of historical interactions, thereby more closely approximating adaptive human behavior in public discussion.

It should also be noted that although the term “reflection” has been systematized mainly in recent LLM-agent research, functionally similar adaptive-feedback ideas already appeared in earlier studies of complex systems and diffusion dynamics. For example, in coupled epidemic-awareness models, individuals adjust their behavior according to risk perception, awareness, or previous exposure, which in turn changes the diffusion process itself~\citep{funk2009awareness, granell2013awareness}. These mechanisms are usually represented by state transitions, risk weights, or behavioral-response parameters rather than natural-language reflection, but they nonetheless embody the same closed-loop logic of using historical experience to revise future decisions. Compared with traditional models, reflection mechanisms in LLM agents extend this loop from numerical state adjustment to semantic-level explanation, emotional regulation, and narrative reconstruction, making them more suitable for modeling reinterpretation, re-judgment, and re-expression in public-emergency opinion dynamics.

However, most existing reflection mechanisms are still designed for single-agent problem solving, embodied decision-making, or general interaction optimization, where the main objective is to improve task success, planning ability, or tool use~\citep{yao2023react, shinn2023reflexion}. In public-emergency opinion simulation, reflection should serve not only the goal of “solving tasks correctly,” but also the reproduction of social cognitive processes, such as how users revise their stance, regulate emotion, or reinforce prior beliefs after receiving official information, reading others' comments, and observing event progression. Motivated by this distinction, this paper incorporates a reflection mechanism into a time-aware multi-agent public-opinion simulation framework to model cognitive updating and expressive adjustment in continuous event streams more faithfully.
社交网络研究范式大致经历了从静态结构分析到动态过程建模的演进，基于代理的建模（agent-based modeling, ABM）逐步成为模拟和分析动态社交网络的重要工具~\citep{li2026abm, gilbert2008agent, wang2025abm}. ABM 框架通常包含代理、环境、交互与规则四个核心要素，通过建模个体的行为模式和交互规则，进而观察个体层面的行为如何经由局部机制演化为宏观社会结构或群体性现象~\citep{gilbert2008agent, vespignani2012modelling, gaoj2016universal}. 早期，ABM 常与传染病模型和意见动力学模型相结合，被用于模拟流行病在复杂社会网络中的传播动态、完全连接网络中的意见演化过程，以及突发公共事件中的恐慌性购买行为~\citep{wu2024epidemic, alatas2017opinion, fan2024panic}. 这类方法在刻画总体扩散速度、网络拓扑效应和宏观演化形态方面具有明显优势；但正如已有研究所指出的，它们通常依赖有限的状态转移规则与逻辑函数，较难在同一系统中同时表达“看见—理解—反应—再表达”的细粒度认知过程，也较难处理现实舆情中持续出现的解释、再叙述与情绪调节现象~\citep{li2026abm, gaoc2024survey}.

\subsection{LLM-Driven Multi-Agent Social Simulation}

大语言模型的发展为多智能体社会仿真提供了新的可能。Park 等~\citep{park2023generative} 提出的生成式智能体利用 LLM 模拟人类角色的日常对话、记忆提取与长期规划，由此推动社交网络研究从“规则驱动的符号世界”向“语义驱动的认知世界”转变~\citep{li2026abm}. 基于 LLM 的多智能体系统能够在角色扮演、文本理解、评论生成和情境推断层面提供更高的语义灵活性，使研究者有机会将个体认知、群体互动与平台环境纳入同一个仿真闭环~\citep{gaoc2024survey, park2023generative}. 近年来，多项研究将这一路径拓展至不同场景：Mou 等~\citep{mou2024socialmovement} 采用“核心 LLM 智能体 + 普通 ABM 智能体”的混合框架，对社会运动中的集体意见变化进行了模拟；Yang 等~\citep{yang2024oasis} 提出的 OASIS 系统实现了百万级智能体的社交媒体仿真；Gao 等~\citep{gao2023s3} 的 S$^3$ 系统和 Piao 等~\citep{piao2025agentsociety} 的 AgentSociety 分别在千人和大规模社会环境中模拟了平台信息传播与集体行为。在谣言与对抗性场景方面，Liu 等~\citep{liu2025rumorsphere} 的 RumorSphere 通过动态交互策略实现了百万级谣言传播仿真，Zhang 等~\citep{zhang2024trendsim} 的 TrendSim 则引入时间感知交互机制以模拟攻击情境下的话题操纵。总体而言，这些工作在技术规模与生成能力上取得了显著进展，但其研究框架仍多围绕异常传播、对抗操纵或特定任务场景展开。对于真实突发公共事件而言，官方发布时间、发布频率、内容立场以及介入强度会在不同阶段对舆情轨迹产生调节作用，而现有模型对这类“官方介入变量”的系统建模仍较为有限。

\subsection{Reflection Mechanisms in LLM Agents}

反思（reflection）是使智能体根据既往经历更新判断与行为策略的重要认知机制。就 LLM 智能体而言，Park 等~\citep{park2023generative} 已将 reflection 作为生成式智能体架构中的关键模块，通过对既往记忆进行归纳，形成较高层次的总结性认知，从而影响后续计划与交互。此后，Yao 等~\citep{yao2023react} 提出的 ReAct 框架证明了将推理轨迹与行动过程交错组织，能够帮助模型在环境交互中持续更新中间判断；Shinn 等~\citep{shinn2023reflexion} 提出的 Reflexion 框架则进一步把语言反馈、自我评估与再尝试结合起来，使智能体能够基于失败经验生成“语言化反思”，并在下一轮决策中显式利用这些反思结果。对于多智能体仿真而言，这类设计的价值不只在于提高任务完成率，更在于使智能体能够依据历史互动对自身认知状态、情绪倾向与表达策略进行动态修正，从而更接近真实个体在公共讨论中的适应性行为。

需要指出的是，尽管“反思”作为术语主要在近年的 LLM 智能体研究中被系统化提出，但在更早期的复杂系统与传播动力学研究中，已经存在与之功能相近的自适应反馈思想。例如，在疾病传播与行为共演化研究中，个体会依据风险感知、信息 awareness 或既往接触经历调整自身行为，从而反过来改变传播过程~\citep{funk2009awareness, granell2013awareness}. 这类机制虽然通常以状态转移、风险权重或行为响应参数的形式出现，而非以自然语言反思的形式实现，但其核心逻辑同样体现了“基于历史经验修正后续决策”的闭环过程。与传统模型相比，LLM 智能体中的反思机制能够把这种闭环从单一数值状态扩展到语义层面的解释、情绪调节与叙事重构，因此更适合刻画突发公共事件舆情中常见的再理解、再判断与再表达过程。

然而，现有研究中的反思机制大多面向单智能体任务求解、具身决策或通用交互优化，其目标通常是提升任务成功率、规划能力或工具使用效果~\citep{yao2023react, shinn2023reflexion}. 在面向公共事件舆情仿真的场景中，反思不仅应服务于“做对任务”，更应服务于社会认知过程的再现，例如用户如何在接触官方信息、他人评论与事件进展后修正立场、缓和情绪或强化既有判断。基于这一差异，本文将反思机制引入时间感知的多智能体舆情仿真框架，以更细致地建模个体在连续事件流中的认知更新与表达调整。
